\subsection{Lifting of paths}

THEOREM 3. --- Let I be an interval of R and let X be a nonempty separated topological space. Let p: X → I be a continuous, open, and proper map whose fibers are totally disconnected (TG, I, p.~83). The map p is surjective. For every point x of X, there exists a continuous section s of p such that s(p(x)) = x.

The set \(p(\mathrm{X})\) is an open, closed (TG, I, p.~72, prop. 1), nonempty subset of I. Since I is connected, we have \(p(\mathrm{X}) = \mathrm{I}\) ; the map \(p\) is therefore surjective.

For every pair \((a,b)\in\mathrm{I}\times\mathrm{I}\) such that \(a\leqslant b\) , denote by \(F_{a,b}\) the set of pairs \((y,z)\in\overline{p}^{-1}(a)\times\overline{p}^{-1}(b)\) such that y and z belong to the same connected component of \(\overline{p}^{-1}([a,b])\) .

Lemma 10. --- Let \(a\) and \(b\) be points of I such that \(a \leqslant b\) .

\begin{enumerate}
\def\labelenumi{\alph{enumi})}
\item
  \(F_{a,b}\) is closed in \(\overline{p}^{-1}(a)\times\overline{p}^{-1}(b)\).
\item
  We have \(\operatorname{pr}_1(\mathrm{F}_{a,b}) = \overline{p}^{-1}(a)\) and \(\operatorname{pr}_2(\mathrm{F}_{a,b}) = \overline{p}^{-1}(b)\) .
\item
  Let \(c \in I\) such that \(b \leqslant c\) . If \((y,z)\) belongs to \(F_{a,b}\) and \((z,t)\) belongs to \(F_{b,c}\) then \((y,t)\) belongs to \(F_{a,c}\) .
\end{enumerate}

The set \(\overline{p}^{-1}([a,b])\) is compact (TG, I, p.~77, prop. 7). Consequently, for a pair \((y,z)\in\overline{p}^{-1}(a)\times\overline{p}^{-1}(b)\) to belong to \(F_{a,b}\) , it is necessary and sufficient that every open and closed subset of \(\overline{p}^{-1}([a,b])\) which contains \(y\) contains \(z\) (TG, II, p.~32, prop. 6).

Let us set \(\mathrm{Y} = \overline{p}^{-1}([a,b])\) . For every subset U of Y that is both open and closed, the set \(((\mathrm{U} \times \mathrm{U}) \cup ((\mathrm{Y} - \mathrm{U}) \times (\mathrm{Y} - \mathrm{U}))) \cap (\overline{p}^{-1}(a) \times \overline{p}^{-1}(b))\) is closed in \(\overline{p}^{-1}(a) \times \overline{p}^{-1}(b)\) . The intersection of these sets is equal to \(\mathrm{F}_{a,b}\) , whence \(a\) ).

Let \(y \in \overline{p}^{-1}(a)\) . Let \(\mathcal{U}\) denote the set of open and closed neighborhoods of y in Y. The map from Y to {[}a, b{]} deduced from p by passing to subspaces is open and proper (I, p.~17, prop. 8), hence also closed. It follows that, for every \(U\) belonging to \(\mathcal{U}\) , \(p(U)\) is a nonempty open and closed subset of {[}a, b{]}; since the interval {[}a, b{]} is connected, one has \(p(U) = [a, b]\) and in particular \(U \cap \overline{p}^{-1}(b) \neq \varnothing\) . Consequently, \((U \cap \overline{p}^{-1}(b))_{U \in \mathcal{U}}\) is a decreasing filtered family of nonempty closed subsets of the compact space \(\overline{p}^{-1}(b)\) . The intersection of this family is nonempty (TG, I, p.~59); let z be one of its elements. One has \((y, z) \in \mathrm{F}_{a,b}\) . We have proved the relation \(\mathrm{pr}_{1}(\mathrm{F}_{a,b}) = \overline{p}^{-1}(a)\) . The relation \(\mathrm{pr}_{2}(\mathrm{F}_{a,b}) = \overline{p}^{-1}(b)\) is deduced from it by replacing p, I, a, b by -p, -I, -b, -a.

Under the hypotheses of \(c\) ), the pair \((y,t)\) belongs to \(\overline{p}^{-1}(a)\times\overline{p}^{-1}(c)\) . The set \(\{y,z\}\) is contained in a connected subset C of \(\overline{p}^{-1}([a,b])\) and the set \(\{z,t\}\) is contained in a connected subset C' of \(\overline{p}^{-1}([b,c])\) . Then, C \(\cup\) C' is a connected subset of \(\overline{p}^{-1}([a,c])\) (TG, I, p.~81, prop. 2) and contains \(\{y,t\}\) , whence the relation \((y,t)\in\mathrm{F}_{a,c}\) .

Returning to the proof of Th. 3. Each fiber of the map p is compact (TG, I, p.~77, prop. 7). By Tychonoff's theorem (TG, I, p.~63, th. 3), the product space \(K = \prod_{a \in I} \overline{p}^{-1}(a)\) is compact. Let \(K'\) be the set of elements \((y_{a})_{a \in I}\) of K such that \(y_{p(x)}\) equals x and that one has \((y_{a}, y_{b}) \in F_{a,b}\) for every pair \((a, b)\) of elements of I such that a \textless{} b. Theorem 3 follows from the following lemma.

Lemma 11. --- a) The set K' is not empty.

\begin{enumerate}
\def\labelenumi{\alph{enumi})}
\setcounter{enumi}{1}
\tightlist
\item
  Let \((s_{a})_{a\in\mathrm{I}}\) be an element of \(K'\) . The map \(s\colon a\mapsto s_{a}\) from I to X is a continuous section of p such that \(s(p(x))=x\) .
\end{enumerate}

For every finite subset S of I containing the point \(p(x)\) , denote by \(K_{S}\) the set of elements \((y_{a})_{a\in\mathrm{I}}\) of K satisfying the relation \(y_{p(x)} = x\) and the relations \((y_{a}, y_{b}) \in \mathrm{F}_{a,b}\) for every pair \((a, b)\) of elements of S such that a \textless{} b. The sets \(K_{S}\) are closed in K (Lemma 10, a)) and form a decreasing filtered family of subsets of K, whose intersection \(K'\) .

To prove that \(K'\) is nonempty, it suffices to show that, for every finite subset S of I containing the point \(p(x)\) , the set \(K_{S}\) is not empty (TG, I, p.~59).

Let S be such a subset; arrange its elements in a strictly increasing sequence \((a_{1},\ldots,a_{n})\), and let i denote the integer such that \(p(x)=a_{i}\) . Set \(y_{a_{i}}=x\) . Lemma 10, b), allows one to construct by induction elements \(y_{a_{j}}\in\overline{p}^{-1}(a_{j})\) , for \(i<j\leqslant n\) , and, by descending induction, elements \(y_{a_{j}}\in\overline{p}^{-1}(a_{j})\) for \(1\leqslant j<i\) , so that one has \((y_{a_{j}},y_{a_{j+1}})\in\mathrm{F}_{a_{j},a_{j+1}}\) for every integer j such that \(1\leqslant j<n\) . By Lemma 10, c), we have \((y_{a},y_{b})\in\mathrm{F}_{a,b}\) for every pair \((a,b)\) of elements of S such that a\textless b. Since the map p is surjective, we can choose for every \(a\in I-S\) an element \(y_{a}\in\overline{p}^{-1}(a)\) . The family \((y_{a})_{a\in I}\) thus constructed belongs to \(K_{S}\) , therefore \(K_{S}\) is not empty.

Let us prove \(b\) ). By definition of \(K'\) , \(s\) is a section of \(p\) such that \(s(p(x)) = x\) . Let \(a \in I\) . Let us prove the continuity of \(s\) at the point \(a\) . Let U be an open neighborhood of \(s_a\) in X. Since \(\overline{p}^{-1}(a)\) is a compact space (TG, I, p.~77, prop. 7) and totally disconnected, \(s_a\) has, in \(\overline{p}^{-1}(a)\) , an open and closed neighborhood C contained in U (TG, II, p.~32, corollary). The sets C and \(\overline{p}^{-1}(a) - C\) are closed in \(\overline{p}^{-1}(a)\) ; therefore compact, and they are disjoint. Since X is separated, they have disjoint open neighborhoods V and \(V'\) . The set \((V \cap U) \cup V'\) is an open neighborhood of the fiber \(\overline{p}^{-1}(a)\) in X; since the map \(p\) is closed (TG, I, p.~72, prop. 1), \((V \cap U) \cup V'\) contains a set of the form \(\overline{p}^{-1}(J)\) , where \(J \subset I\) is an open interval containing \(a\) (I, p.~75, lemma). Put \(W = V \cap U \cap \overline{p}^{-1}(J)\) . The set W is open in \(\overline{p}^{-1}(J)\) ; it is also closed in \(\overline{p}^{-1}(J)\) since we have \(\overline{p}^{-1}(J) - W = V' \cap \overline{p}^{-1}(J)\) . Let \(b \in J\) . The closed interval of I with endpoints \(a\) and \(b\) is contained in J. By hypothesis, \((s_a, s_b)\) belongs to \(F_{a,b}\) if \(a \leqslant b\) and \((s_b, s_a)\) belongs to \(F_{b,a}\) if \(b \leqslant a\) . There therefore exists a connected subset of \(\overline{p}^{-1}(J)\) containing \(\{s_a, s_b\}\) . Consequently, the point \(s_b\) belongs to every open and closed subset of \(\overline{p}^{-1}(J)\) which contains \(s_a\) , hence in particular to W; a fortiori, \(s_b\) belongs to U. We therefore have \(s(J) \subset U\) , which proves the continuity of \(s\) at the point \(a\) .

COROLLARY. --- Let X and B be topological spaces and \(p: X \to B\) a continuous, open, proper and separated map whose fibers are totally disconnected. Let I be an interval of R, let \(f: I \to B\) be a continuous map, let a be a point of I and let x be a point of X such that \(f(a) = p(x)\) . There exists a continuous map \(g: I \to X\) such that \(p \circ g = f\) and \(g(a) = x\) .

Let us set \(X' = I \times_{B} X\) and denote by \(p'\colon X'\to I\) and \(f'\colon X'\to X\) the canonical projections. The map \(p'\) is continuous, open, proper and separated (I, p.~17, prop. 8 and p.~27, prop. 4). Since the space I is separated, the space \(X'\) is separated (I, p.~26, remark 3). The fibers of \(p'\) are totally disconnected (I, p.~10, corollary, a)). By th. 3, there exists a continuous section \(s'\) of \(p'\) which takes the value \((a,x)\) at \(a\) . The map \(g = f'\circ s'\) from I to X is continuous; we have \(p\circ g = p\circ f'\circ s' = f\circ p'\circ s' = f\) and \(g(a) = x\) , whence the corollary.

THEOREM 4. --- Let X be a topological space, let G be a discrete group acting properly on X and let p: X → X/G be the canonical map. Let I be an interval of R, let f: I → X/G be a continuous map, let a be a point of I and let x be a point of X such that \(f(a) = p(x)\) . There exists a continuous map \(\varphi\) : I → X such that \(p \circ \varphi = f\) and \(\varphi(a) = x\) .

Let us first treat the case where I is a closed bounded interval of R.

By TG, III, p.~29, prop. 3, the space X is separated.

Let \(y\) be a point of X/G and let \(x \in X\) such that \(y = p(x)\) . The stabilizer \(K_x\) of \(x\) is a finite subgroup of G; moreover, there exist neighborhoods \(U_x\) of \(x\) in X and \(V_y\) of \(y\) in X/G such that \(U_x\) is stable under \(K_x\) , \(gU_x \cap U_x = \varnothing\) if \(g \notin K_x\) , and \(p\) induces a homeomorphism from \(U_x/K_x\) onto \(V_y\) . In addition, for every \(g \in G\) , \(p\) induces a homeomorphism from \(gU_x\) onto \(V_y\) . Since I is compact, there exist integers \(m\) and \(n\) such that \(m \leqslant 0 \leqslant n\) and a finite sequence \((a_i)_{m \leqslant i \leqslant n}\) of elements of I such that \(a_0 = a\) , \(I = {[}a_m,a_n{]}\) , and such that, for every \(i \in \{m, \ldots, n-1\}\) , \(f({[}a_i,a_{i+1}{]})\) is contained in an open set \(V_{y_i}\) of X/G constructed as above.

Let \(x_0\) be the unique element of \(\overline{p}^{-1}(y_0)\) such that \(x \in \mathrm{U}_{x_0}\) . Let \(q_0\) be the canonical map from \(\mathrm{U}_{x_0}\) onto \(\mathrm{U}_{x_0}/\mathrm{K}_{x_0}\) by passing to the quotient. The map \(p\) induces a homeomorphism \(i_0\) of \(\mathrm{U}_{x_0}/\mathrm{K}_{x_0}\) onto \(\mathrm{V}_{y_0}\) such that \(i_0 \circ q_0 = p \mid \mathrm{U}_{x_0}\) . The map \(q_0\) is proper (TG, III, p.~29, prop. 3), open (TG, I, p.~31, example 1), and separated, since X is separated. Its fibers are totally disconnected, since they are finite. By the corollary, III, p.~286, there exists a continuous map \(\varphi_0\colon [a_0,a_1]\to \mathrm{U}_{x_0}\) such that \(p\circ \varphi_0 = f\mid [a_0,a_1]\) .

We construct similarly, by induction on the integer \(i \in \{0, \ldots, n-1\}\) , a point \(x_i \in \overline{p}^{-1}(y_i)\) and a continuous map \(\varphi_i: [a_i, a_{i+1}] \to X\) whose image is contained in \(\mathrm{U}_{x_i}\) such that \(p \circ \varphi_i = f \mid [a_i, a_{i+1}]\) and such that \(\varphi_i(a_{i+1}) = \varphi_{i+1}(a_{i+1})\) if \(0 \leqslant i < n-1\) .

Analogously, one constructs by descending induction on the integer \(i \in \{m, \ldots, -1\}\) , a point \(x_i \in \overline{p}^{-1}(y_i)\) and a continuous map \(\varphi_i: [a_i, a_{i+1}] \to X\) whose image is contained in \(U_{x_i}\) such that \(p \circ \varphi_i = f \mid [a_i, a_{i+1}]\) and such that \(\varphi_i(a_{i+1}) = \varphi_{i+1}(a_{i+1})\) if \(m \leqslant i < 0\) .

There exists a unique map \(\varphi\colon\mathrm{I}\to\mathrm{X}\) which coincides with \(\varphi_{i}\) on \([a_{i},a_{i+1}]\) for \(m\leqslant i<n\) . It is continuous (TG, I, p.~19, prop. 4). It is a continuous lifting of \(f\) to X such that \(\varphi(a)=x\) .

This proves the theorem when I is compact. In the general case, there exist sequences \((a_{n})_{n\in\mathbf{N}}\) and \((b_{n})_{n\in\mathbf{N}}\) such that \((a_{n})\) is stationary with limit \(\inf(I)\), \((b_{n})\) is stationary with limit \(\sup(I)\), \(a=a_{0}=b_{0}\) , and such that \((a_{n})\) (resp. \((b_{n})\) ) is constant if I has a least (resp. a greatest) element. By the preceding, for every integer \(n\in N\) there exists a continuous lifting \(\varphi_{n}\) of \(f\mid[a_{n},a_{n+1}]\) to X and a continuous lifting \(\varphi_{n}^{\prime}\) of \(f\mid[b_{n+1},b_n]\) to X such that \(\varphi_{0}(a_{0})=\varphi_{0}^{\prime}(b_{0})=x\) and \(\varphi_{n+1}(a_{n+1})=\varphi_{n}(a_{n+1})\), \(\varphi_{n+1}^{\prime}(b_{n+1})=\varphi_{n}^{\prime}(b_{n+1})\) . There exists a unique map \(\varphi\colon I\to X\) which coincides with \(\varphi_{n}\) on \([a_{n},a_{n+1}]\) and with \(\varphi_{n}^{\prime}\) on \([b_{n+1},b_{n}]\) , for every \(n\in N\) . It is continuous (loc. cit.) and it is a continuous lifting of f to X such that \(\varphi(a)=x\) . The theorem is thus proved.

Examples. --- 1) Let X be a separated topological space and let G be a finite group, endowed with the discrete topology, acting continuously on X. The action is then proper (TG, III, p.~28, prop. 2). The assertion of Theorem 4 then follows directly from the corollary of Theorem 3.

\begin{enumerate}
\def\labelenumi{\arabic{enumi})}
\setcounter{enumi}{1}
\tightlist
\item
  Let n be an integer \(\geqslant 0\) . Denote by \(P_{n}\) the set of monic polynomials \(P \in C[X]\) of degree n, endowed with the topology for which the map \((c_{0}, \ldots, c_{n-1}) \mapsto X^{n} + c_{n-1}X^{n-1} + \cdots + c_{0}\) is a homeomorphism of \(C^{n}\) onto \(P_{n}\) . The map p from \(C^{n}\) to \(P_{n}\) defined by \(p(z_{1}, \ldots, z_{n}) = (X - z_{1}) \ldots (X - z_{n})\) is continuous. The symmetric group \(S_{n}\) acts on \(C^{n}\) by permutation of the factors and p defines, by passage to the quotient, a homeomorphism from \(C^{n}/S_{n}\) onto \(P_{n}\) (TG, VIII, p.~22, prop. 1, I, p.~23, cor. 1 and TG, VIII, p.~20). We therefore deduce the following statement:
\end{enumerate}

Let I be an interval of R, let \((c_{0},\ldots,c_{n-1})\) be a sequence of continuous maps from I to C, let a be a point of I and let \((z_{1},\ldots,z_{n})\) be a sequence of complex numbers such that \((\mathrm{X}-z_{1})\ldots(\mathrm{X}-z_{n})=\mathrm{X}^{n}+c_{n-1}(a)\mathrm{X}^{n-1}+\cdots+c_{0}(a)\) . There exists a sequence \((\lambda_{1},\ldots,\lambda_{n})\) of continuous maps from I to C such that \(\lambda_{i}(a)=z_{i}\) for \(1\leqslant i\leqslant n\) and \((\mathrm{X}-\lambda_{1}(t))\ldots(\mathrm{X}-\lambda_{n}(t))=\mathrm{X}^{n}+c_{n-1}(t)\mathrm{X}^{n-1}+\cdots+c_{0}(t)\) for all \(t\in I\) .

